\chapter{Problem And Scope}

\section{Project Goal}

\texttt{apfs-fastindex} is investigating whether an APFS indexer can eventually
deliver WizTree-like scans by reading APFS metadata structures directly instead
of walking every file through normal filesystem APIs.

The narrow v1 goal is smaller:

\begin{itemize}
  \item one APFS volume
  \item one coherent selected state
  \item namespace reconstruction
  \item logical size
  \item fail-closed behavior outside the tested allowlist
\end{itemize}

The canonical product and v1 scope sources are \path{spec.md} and
\path{docs/research/contracts/narrow-v1-parser-contract.md}.
The long-range route from the narrow parser to a hybrid consumer product is
tracked in
\path{docs/research/plans/general-wiztree-for-any-mac-roadmap.md}, which
defines staged Gates 0--8 from the narrow Rust full scan to a hybrid product
and a parallel R0--R4 release shape.

\section{Why APFS Is Not NTFS MFT Walking}

The project cannot simply copy an NTFS MFT strategy. APFS uses copy-on-write
B-trees, object maps, checkpointed transaction state, volume roles, snapshots,
clones, sparse files, compression, and modern macOS presentation layers. The
research therefore starts with correctness and state selection before speed.

\section{Current Non-Goals}

The current raw v1 path does not claim:

\begin{itemize}
  \item live startup-disk raw scanning
  \item merged System/Data/Finder-visible root semantics
  \item physical, shared, exclusive, or snapshot-retained accounting
  \item persistent incremental cache reuse
  \item production performance
\end{itemize}

Those remain research tracks, not implementation promises.
